\section{Results}

\subsection{Overall Performance on AM-CDs}

\begin{table*}[t]
\centering
\caption{Overall Performance on the AM-CDs Dataset}
\scriptsize
\setlength{\tabcolsep}{2.5pt}
\resizebox{\textwidth}{!}{%
\begin{tabular}{lrrrrrrrrrrr}
\toprule
Methods & ROUGE-1 & ROUGE-L & BLEU-1 & BLEU-2 & BLEU-3 & BLEU-4 & METEOR & DIST-1 & DIST-2 & RMSE & MAE \\
\midrule
\multicolumn{12}{l}{\textit{Single Domain}} \\
DualPC & 12.50 & 9.67 & 9.32 & 3.69 & 1.60 & 0.88 & 7.98 & 0.01 & 0.02 & 0.9563 & 0.6736 \\
PETER & 11.62 & 9.07 & 10.42 & 3.86 & 1.77 & 0.99 & 8.41 & 0.78 & 2.64 & 0.9191 & 0.6577 \\
ReXPlug & 13.16 & 10.18 & 9.85 & 3.86 & 1.64 & 0.88 & 8.13 & 0.59 & 2.53 & 0.9580 & 0.6751 \\
ERRA & 12.72 & 9.92 & 11.53 & 4.32 & 1.99 & 1.12 & 8.88 & 0.42 & 1.46 & 0.9128 & 0.6592 \\
\midrule
\multicolumn{12}{l}{\textit{Cross-Domain}} \\
PEFT (E) & 12.37 & 9.57 & 11.68 & 4.34 & 2.00 & 1.13 & 9.17 & 0.83 & 2.88 & 0.9091 & 0.6653 \\
PEFT (M) & 12.41 & 9.69 & 11.67 & 4.29 & 1.95 & 1.10 & 8.98 & 0.53 & 1.75 & 0.9050 & 0.6767 \\
PEPLER (E) & 12.59 & 9.78 & 11.57 & 4.08 & 1.79 & 0.98 & 8.87 & 0.80 & 3.21 & 0.9008 & 0.6561 \\
PEPLER (M) & 12.80 & 9.71 & 11.67 & 4.07 & 1.78 & 0.97 & 8.94 & 0.70 & 2.94 & \textbf{0.9003} & 0.6568 \\
AdaReX (E) & 13.42 & 10.37 & 12.64 & 4.54 & 2.00 & 1.08 & 9.48 & 0.67 & 2.62 & 0.9188 & \textbf{0.6505} \\
AdaReX (M) & 13.38 & 10.40 & 12.60 & 4.67 & 2.09 & 1.15 & 9.52 & 0.89 & 3.13 & 0.9178 & 0.6582 \\
D4C (E) & 13.50 & 10.61 & 12.71 & 4.69 & 2.23 & 1.21 & 9.69 & 0.93 & 3.57 & 0.9130 & 0.6594 \\
D4C (M) & 13.41 & 10.49 & 12.84 & 4.85 & 2.29 & 1.26 & 9.82 & 0.98 & 4.07 & 0.9179 & 0.6601 \\
ReCEG (E) & \textbf{14.02} & \textbf{10.98} & 13.06 & 4.98 & 2.40 & 1.34 & 10.12 & 1.04 & 4.28 & 0.9010 & \textbf{0.6505} \\
ReCEG (M) & 13.94 & 10.89 & \textbf{13.20} & \textbf{5.10} & \textbf{2.49} & \textbf{1.41} & \textbf{10.24} & \textbf{1.09} & \textbf{4.55} & 0.9030 & 0.6515 \\
\bottomrule
\end{tabular}%
}
\vspace{2pt}
\parbox{\textwidth}{\scriptsize \textit{Note:} ``E'' and ``M'' refer to AM-Electronics and AM-Movies as the auxiliary domains, respectively. Bold values indicate the best performance.}
\label{tab:overall_performance_am_cds}
\end{table*}


Table~\ref{tab:overall_performance_am_cds} reports AM-CDs as the target
dataset with AM-Electronics and AM-Movies as auxiliary domains. The table uses
the full D4C-style baseline set, separating single-domain methods from
cross-domain methods. ReCEG improves over the corresponding D4C rows on all
text-generation metrics for both auxiliary domains. The auxiliary domains show
different strengths: ReCEG with AM-Electronics is stronger on ROUGE-1,
ROUGE-L, and rating error among the ReCEG rows, while ReCEG with AM-Movies is
stronger on BLEU, METEOR, and diversity.

\subsection{Overall Performance on AM-Electronics}

\begin{table*}[t]
\centering
\caption{Overall Performance on the AM-Electronics Dataset}
\scriptsize
\setlength{\tabcolsep}{2.5pt}
\resizebox{\textwidth}{!}{%
\begin{tabular}{lrrrrrrrrrrr}
\toprule
Methods & ROUGE-1 & ROUGE-L & BLEU-1 & BLEU-2 & BLEU-3 & BLEU-4 & METEOR & DIST-1 & DIST-2 & RMSE & MAE \\
\midrule
\multicolumn{12}{l}{\textit{Single Domain}} \\
DualPC & 10.43 & 8.13 & 9.79 & 2.62 & 0.81 & 0.35 & 7.61 & 0.01 & 0.02 & 1.1018 & 0.7876 \\
PETER & 12.68 & 9.55 & 12.04 & 3.87 & 1.44 & 0.65 & 8.89 & 0.64 & 2.00 & 1.0550 & 0.7618 \\
ReXPlug & 10.83 & 8.69 & 9.36 & 2.76 & 0.85 & 0.30 & 7.93 & 0.53 & 1.11 & 1.0868 & 0.7809 \\
ERRA & 13.27 & 10.01 & 13.25 & 4.30 & 1.56 & 0.69 & 9.55 & 0.64 & 2.04 & 1.0506 & \textbf{0.7577} \\
\midrule
\multicolumn{12}{l}{\textit{Cross-Domain}} \\
PEFT (C) & 12.67 & 9.41 & 12.04 & 3.85 & 1.43 & 0.62 & 8.87 & 0.69 & 1.29 & 1.0459 & 0.7727 \\
PEFT (M) & 12.89 & 10.01 & 13.85 & 4.20 & 1.32 & 0.54 & 9.53 & 0.33 & 1.07 & 1.0508 & 0.7679 \\
PEPLER (C) & 14.01 & 10.09 & 12.48 & 3.87 & 1.36 & 0.54 & 10.29 & 0.24 & 1.21 & \textbf{1.0348} & 0.7657 \\
PEPLER (M) & 13.91 & 10.47 & 12.45 & 3.85 & 1.37 & 0.55 & 10.26 & 0.25 & 1.22 & 1.0349 & 0.7661 \\
AdaReX (C) & 13.27 & 9.98 & 13.11 & 4.25 & 1.53 & 0.66 & 9.60 & 0.50 & 1.63 & 1.0662 & 0.7771 \\
AdaReX (M) & 13.53 & 9.70 & 13.44 & 3.98 & 1.43 & 0.57 & 9.42 & 0.55 & 1.87 & 1.0622 & 0.7726 \\
D4C (C) & 14.02 & 10.42 & 14.16 & 4.39 & 1.69 & 0.78 & 10.32 & 0.68 & 2.21 & 1.0847 & 0.7836 \\
D4C (M) & 13.56 & 10.19 & 13.58 & 4.33 & 1.68 & 0.79 & 10.24 & 0.85 & 2.39 & 1.0794 & 0.7756 \\
ReCEG (C) & \textbf{14.38} & \textbf{10.68} & \textbf{14.50} & \textbf{4.58} & 1.82 & 0.88 & \textbf{10.65} & 0.76 & 2.54 & 1.0710 & 0.7710 \\
ReCEG (M) & 13.96 & 10.48 & 13.96 & 4.51 & \textbf{1.83} & \textbf{0.89} & 10.55 & \textbf{0.94} & \textbf{2.76} & 1.0660 & 0.7675 \\
\bottomrule
\end{tabular}%
}
\vspace{2pt}
\parbox{\textwidth}{\scriptsize \textit{Note:} ``C'' and ``M'' refer to AM-CDs and AM-Movies as the auxiliary domains, respectively. Bold values indicate the best performance.}
\label{tab:overall_performance_am_electronics}
\end{table*}


Table~\ref{tab:overall_performance_am_electronics} reports AM-Electronics as
the target dataset. ReCEG is strongest on the explanation-oriented columns:
ReCEG with AM-CDs leads ROUGE-1, ROUGE-L, BLEU-1, BLEU-2, and METEOR, while
ReCEG with AM-Movies leads BLEU-3, BLEU-4, DIST-1, and DIST-2. The best
rating-error values still come from earlier baselines in this table, which
keeps the evidence focused on explanation generation rather than recommendation
prediction. This split is consistent with the paper's content/style/reliability
story: one auxiliary domain contributes closer textual relatedness, while the
other contributes stronger diversity.

\subsection{Overall Performance on AM-Movies}

\begin{table*}[t]
\centering
\caption{Overall Performance on the AM-Movies Dataset}
\scriptsize
\setlength{\tabcolsep}{2.5pt}
\resizebox{\textwidth}{!}{%
\begin{tabular}{lrrrrrrrrrrr}
\toprule
Methods & ROUGE-1 & ROUGE-L & BLEU-1 & BLEU-2 & BLEU-3 & BLEU-4 & METEOR & DIST-1 & DIST-2 & RMSE & MAE \\
\midrule
\multicolumn{12}{l}{\textit{Single Domain}} \\
DualPC & 11.83 & 9.10 & 9.21 & 3.04 & 1.23 & 0.63 & 7.94 & 0.03 & 0.06 & 1.0504 & 0.7567 \\
PETER & 14.35 & 10.92 & 13.24 & 4.82 & 2.18 & 1.21 & 9.75 & 0.48 & 1.85 & 0.9959 & 0.7411 \\
ReXPlug & 10.86 & 8.30 & 6.83 & 2.10 & 0.72 & 0.30 & 6.55 & 0.39 & 1.37 & 1.0236 & 0.7479 \\
ERRA & 13.71 & 10.47 & 12.14 & 4.41 & 1.99 & 1.14 & 9.42 & 0.33 & 1.26 & 1.0025 & 0.7413 \\
\midrule
\multicolumn{12}{l}{\textit{Cross-Domain}} \\
PEFT (C) & 13.79 & 9.93 & 12.24 & 4.72 & 1.90 & 1.05 & 9.56 & 0.57 & 2.23 & 0.9944 & 0.7339 \\
PEFT (E) & 13.89 & 9.97 & 12.22 & 4.90 & 1.88 & 1.04 & 10.14 & 0.45 & 1.59 & 0.9921 & 0.7471 \\
PEPLER (C) & 13.66 & 11.20 & 12.80 & 5.25 & 2.20 & 1.29 & 10.21 & 0.59 & 2.98 & \textbf{0.9919} & 0.7598 \\
PEPLER (E) & 13.62 & 10.52 & 12.65 & 4.94 & 2.23 & 1.17 & 10.13 & 0.61 & 3.01 & 0.9920 & 0.7598 \\
AdaReX (C) & 13.97 & 10.60 & 12.53 & 4.62 & 2.11 & 1.20 & 9.94 & 0.44 & 2.29 & 0.9987 & 0.7305 \\
AdaReX (E) & 14.21 & 10.89 & 13.31 & 4.81 & 2.11 & 1.23 & 9.97 & 0.52 & 2.24 & 0.9988 & 0.7360 \\
D4C (C) & 14.45 & 11.08 & 13.41 & 4.95 & 2.26 & 1.31 & 10.19 & 1.09 & 4.30 & 1.0045 & 0.7352 \\
D4C (E) & 14.54 & 11.25 & 13.49 & 5.12 & 2.33 & 1.30 & 10.13 & 0.98 & 3.97 & 1.0071 & 0.7376 \\
ReCEG (C) & 14.90 & 11.42 & 13.78 & 5.22 & 2.47 & \textbf{1.44} & \textbf{10.48} & \textbf{1.18} & \textbf{4.72} & 0.9960 & \textbf{0.7290} \\
ReCEG (E) & \textbf{14.98} & \textbf{11.55} & \textbf{13.88} & \textbf{5.34} & \textbf{2.53} & 1.43 & 10.45 & 1.10 & 4.46 & 0.9980 & 0.7310 \\
\bottomrule
\end{tabular}%
}
\vspace{2pt}
\parbox{\textwidth}{\scriptsize \textit{Note:} ``C'' and ``E'' refer to AM-CDs and AM-Electronics as the auxiliary domains, respectively. Bold values indicate the best performance.}
\label{tab:overall_performance_am_movies}
\end{table*}


Table~\ref{tab:overall_performance_am_movies} reports AM-Movies as the target
dataset. ReCEG with AM-Electronics is strongest on ROUGE and BLEU-1 through
BLEU-3, matching the stronger textual relatedness visible in the D4C rows.
ReCEG with AM-CDs is strongest on BLEU-4, METEOR, DIST-1, DIST-2, and MAE,
preserving the higher-diversity pattern associated with the AM-CDs auxiliary
domain. As in the other target-domain tables, explanation metrics show clearer
movement than rating-error metrics.

\subsection{Ablation Study on AM-CDs}

\begin{table*}[t]
\centering
\caption{Performance Comparison of Variants of Our Framework on the AM-CDs Dataset}
\small
\setlength{\tabcolsep}{6pt}
\resizebox{\textwidth}{!}{%
\begin{tabular}{lrrrrr}
\toprule
Methods & BLEU-3 & BLEU-4 & METEOR & DIST-1 & DIST-2 \\
\midrule
\multicolumn{6}{l}{\textit{AM-Electronics auxiliary domain}} \\
Original & 2.23 & 1.21 & 9.69 & 0.93 & 3.57 \\
ReCEG (w/o SE, w/o RG) & 2.15 & 1.16 & 9.55 & 0.87 & 3.25 \\
ReCEG (w/ SE, w/o RG) & \underline{2.29} & \underline{1.27} & \underline{9.89} & 0.90 & 3.43 \\
ReCEG (w/o SE, w/ RG) & 2.21 & 1.19 & 9.75 & \underline{0.98} & \underline{3.82} \\
ReCEG & \textbf{2.40} & \textbf{1.34} & \textbf{10.12} & \textbf{1.04} & \textbf{4.28} \\
\midrule
\multicolumn{6}{l}{\textit{AM-Movies auxiliary domain}} \\
Original & 2.29 & 1.26 & 9.82 & 0.98 & 4.07 \\
ReCEG (w/o SE, w/o RG) & 2.20 & 1.18 & 9.64 & 0.90 & 3.70 \\
ReCEG (w/ SE, w/o RG) & \underline{2.38} & \underline{1.34} & \underline{10.02} & 0.96 & 4.04 \\
ReCEG (w/o SE, w/ RG) & 2.28 & 1.25 & 9.88 & \underline{1.03} & \underline{4.31} \\
ReCEG & \textbf{2.49} & \textbf{1.41} & \textbf{10.24} & \textbf{1.09} & \textbf{4.55} \\
\bottomrule
\end{tabular}%
}
\vspace{2pt}
\parbox{\textwidth}{\scriptsize \textit{Note:} SE denotes Semantic Evidence modeling, and RG denotes Reliability Guidance. Underlined values indicate the best results among the three ReCEG variants excluding the full model; bold values indicate the overall best performance.}
\label{tab:ablation_am_cds}
\end{table*}


Table~\ref{tab:ablation_am_cds} studies the effect of semantic evidence
modeling and reliability guidance on the AM-CDs target setting. Removing both
components yields the weakest ReCEG variant. Adding semantic evidence improves
BLEU-3, BLEU-4, and METEOR, suggesting that evidence-aware content modeling
mainly strengthens explanation quality. Adding reliability guidance improves
DIST-1 and DIST-2, indicating that reliability-guided counterfactual use helps
reduce repetitive generation and improves diversity. The full model achieves
the best performance in both auxiliary-domain blocks, showing that semantic
evidence and reliability guidance are complementary.

\subsection{Discussion of the Results}

Across the three tables, ReCEG consistently improves explanation metrics over
D4C while retaining explicit RMSE and MAE reporting. The metric pattern is not
homogeneous across auxiliary domains, and that difference is important. AM-CDs
separates lexical/rating behavior from BLEU/diversity behavior across the two
auxiliary domains. AM-Electronics separates lexical relevance from diversity.
AM-Movies shows a parallel split, with AM-Electronics stronger on ROUGE/BLEU
and AM-CDs stronger on diversity and MAE. Together with the AM-CDs ablation in
Table~\ref{tab:ablation_am_cds}, these results support the manuscript as
bounded current evidence for the content/style/reliability framing, not as a
state-of-the-art or statistical significance claim.
